% Sandia National Laboratories is a multimission laboratory managed and
% operated by National Technology & Engineering Solutions of Sandia, LLC, a
% wholly owned subsidiary of Honeywell International Inc., for the U.S.
% Department of Energy’s National Nuclear Security Administration under
% contract DE-NA0003525.

% Copyright 2002-2024 National Technology & Engineering Solutions of Sandia,
% LLC (NTESS).

\documentclass[letterpaper]{scrartcl}
\usepackage[hyperindex=true, colorlinks=false]{hyperref}
\usepackage{ltxtable, multirow}
\usepackage{Xyce}
\usepackage{geometry}

\pdfcatalog {/PageMode /UseNone}
\renewcommand{\arraystretch}{1.2}

% Sets the page margins to be the same as the Guides (SAND reports)
\geometry{pdftex, inner=1in, textwidth=6.5in, textheight=9in}

% Gets rid of Section numbers
\setcounter{secnumdepth}{0}

% Set this here once, and use \XyceVersionVar{} in the document
\XyceVersion{7.7}

% ---------------------------------------------------------------------------- %
%
% Set the title, author, and date
%
\title{\XyceTitle{} Parallel Electronic Simulator\\
Version \XyceVersionVar{} Release Notes}

\author{ Sandia National Laboratories}

\date{\today}

% Approved April TBD 2023
% SAND Number SAND2023-TBD

% ---------------------------------------------------------------------------- %
% Start the document

\begin{document}
\maketitle

The \XyceTM{} Parallel Electronic Simulator has been written to support the
simulation needs of Sandia National Laboratories' electrical designers.
\XyceTM{} is a SPICE-compatible simulator with the ability to solve extremely
large circuit problems on large-scale parallel computing platforms, but also
includes support for most popular parallel and serial computers.

For up-to-date information not available at the time these notes were produced,
please visit the \XyceTM{} web page at
{\color{XyceDeepRed}\url{http://xyce.sandia.gov}}.

\tableofcontents
\vspace*{\fill}
\parbox{\textwidth}
{
  \raisebox{0.13in}{\includegraphics[height=0.5in]{snllineblubrd}}
  \hfill
  \includegraphics[width=1.5in]{xyce_flat_white}
}


\newpage

\section*{Highlights}
\begin{XyceItemize}
\item Xyce now handles local variation with the same netlist convention as
      HSPICE and other simulators.  This allows users to more easily simulate
      things like mismatch using commercial PDKs.

\item The parser, expression library and device package have undergone various
      efficiency improvements with respect to parameter handling.  In the case
      of the parser, this speeds up parsing in many cases.  For the device
      package, it speeds up analysis.  In the case of the expression library, it
      often speeds up both parsing and analysis.
\end{XyceItemize}

\section{New Features and Enhancements}

\subsection*{New Devices and Device Model Improvements}
\begin{XyceItemize}
\item The Xyce diode now supports sidewall effects supported by other
  simulators such as ngspice and HSPICE.
\item The open source version of Xyce now includes the published
  version of EKV MOSFET model version 2.6 as the level 260 MOSFET.
  Previously, this model was only available in our ``non free''
  binaries.
\end{XyceItemize}

\subsection*{Enhanced Solver Stability, Performance and Features}
\begin{XyceItemize}

\item The expression library \texttt{TABLE} and \texttt{SPLINE} operators now
     support a ``downsampling'' feature, in which the user can request that a
     large data file can automatically have the number of points used by the
     source reduced.  For \texttt{TABLE} based sources, which are based on a
     PWL linear approximation, this can dramatically speed up simulations, as
     the number of required breakpoints is reduced.

\item The random operators in the expression library (such as \texttt{AGAUSS}
     and \texttt{AUNIF}) now support implicit local and global variation,
     consistent with other simulators.

\item Parameter handling in the parser has been refactored. As a result, the
     parsing and setup time for large process design kits (PDKs) has been
     improved.

\end{XyceItemize}

\subsection*{Interface Improvements}
\begin{XyceItemize}
\item Subcircuit multipliers (\texttt{M=}) previously had to have a fixed value.
  Now, they can be treated like any other parameter, and can change during a simulation
  using commands such as \texttt{.STEP}.

\item Formatting errors in unused subcircuit definitions (\texttt{.SUBCKT}) will no
  longer emit errors, only warnings.  This is limited to issues in the (\texttt{.SUBCKT})
  line of the netlist.  If additional parsing issues are found within the subcircuit
  definition, they may still result in an error being emitted from the simulator.

\item \Xyce{} now has a diagnostic option, \texttt{.options diagnostic} to aid in
  debugging of user circuits.  This option can be used to output nodes with extreme
  solution values, voltages, currents and discontinuities during a simulation.
  Additionally, it can provide information on continuation methods used during
  the DC Operating point calculation.  See the Reference Guide section 2.1.25.2
  and chapter 10 the User's Guide for more details.

\item The third argument (\texttt{num\_sigma}) to the random expression
     operators \texttt{AGAUSS} and \texttt{GAUSS} is now optional.

\item The metrics tracking and reporting feature now has an option to include
  error messages. When enabled, the tracking data reported by Xyce is encoded as
  a JSON string. If certain error conditions occur, the JSON string may now
  include a list of diagnostic messages. These error messages will be encoded as
  an array under the tag name “ErrorMessages” next to the “audit” field
  (e.g. {“audit”:\{\dots\},”ErrorMessages”:[\dots]}).

\item Expressions are now allowed on the following analysis lines:
     \texttt{.TRAN}, \texttt{.DC} and \texttt{.AC}.

\item Data files read in via the expression library operators such as
     \texttt{TABLE} and \texttt{TABLEFILE} can now contain comments.

\item Data files read in via the expression library operators such as
     \texttt{TABLE} and \texttt{TABLEFILE} can now include dashes in the
     filename.

\item When running a frequency domain analysis such as \texttt{.HB} or
     \texttt{.AC}, complex-valued expressions on the \texttt{.PRINT} line will
     now automatically output both the real and imaginary parts.

\item \texttt{.params} can now be complex-valued.  When used in a real-valued
     context, only the real part will be used, but in complex (i.e. frequency
     domain) contexts, complex arithmetic is performed.

\end{XyceItemize}


\subsection*{XDM}
\begin{XyceItemize}
\item Fixed bug where an instance of a SPECTRE device is not translated as an
    instance or a comment if it appears after a definition of a SPECTRE device.
\end{XyceItemize}


\newpage
\section{Defects Fixed in this Release}
% Sandia National Laboratories is a multimission laboratory managed and
% operated by National Technology & Engineering Solutions of Sandia, LLC, a
% wholly owned subsidiary of Honeywell International Inc., for the U.S.
% Department of Energy's National Nuclear Security Administration under
% contract DE-NA0003525.

% Copyright 2002-2025 National Technology & Engineering Solutions of Sandia,
% LLC (NTESS).

%%
%% Fixed Defects.
%%
{
\small

\begin{longtable}[h] {>{\raggedright\small}m{2in}|>{\raggedright\let\\\tabularnewline\small}m{3.5in}}
    \caption{Fixed Defects.  The Xyce team has multiple issue
     trackers, and the table below indicates fixed issues by
     indentifying both the tracker and the issue number.  Further,
     some issues are reported by open source users on GitHub and these
     issues may be tracked using multiple issue numbers.} \\ \hline
     \rowcolor{XyceDarkBlue} \color{white}\textbf{Defect} & \color{white}\textbf{Description} \\ \hline
     \endfirsthead
     \caption[]{Fixed Defects.  Note that we have two multiple issue tracking systems for Sandia Users.
     SON and SRN refer to our legacy open- and restricted-network Bugzilla system, and Gitlab refers to issues in our gitlab repositories.  } \\ \hline
     \rowcolor{XyceDarkBlue} \color{white}\textbf{Defect} & \color{white}\textbf{Description} \\ \hline
     \endhead

  \textbf{Xyce Project Backlog/888}: Xyce raw files (if specified via the command line "-r" argument) can incorrectly include adjoint sentivity text
  &  There was a logic error in the output manager of \Xyce{}.  This has been fixed.
  \\\hline

  \textbf{Xyce Project Backlog/849}: Sensitivities don't work with PWL source & Analytical sensitivity functions were not implemented for independent sources.  Also PWL sources had unique parsing issues that prevented any kind of derivative from working.  These issues have been addressed.
  \\\hline

  \textbf{Xyce Project Backlog/979}: Fix incorrect output of FD ASCII rawfile in sweeps & There was a logic error in the output manager of \Xyce{} that overwrites the number of points in a Raw ASCII file for AC analysis.  This has been fixed.
  \\\hline
  
  \textbf{Xyce Project Backlog/988}: Fix the capacitance issue in the BSIM SOI 4.6.1 &  There was an error in one of the capacitance models of the BSIM SOI version 4.6.1, resulting in a capacitance that was too large.  This was not fixed in the official version of the 4.6.1 from the BSIM group, but was fixed in the 4.6.1 implementation of commercial simulators.  As such, we chose to fix it in our own implementation. 
  \\\hline

\textbf{Xyce Project Backlog/xxx}: Description
  &  Details
  \\\hline


\end{longtable}
}


\newpage
\section{Supported Platforms}
\subsection*{Certified Support}
The following platforms have been subject to certification testing for the
\Xyce{} version 7.7 release.
\begin{XyceItemize}
  \item Red Hat Enterprise Linux${}^{\mbox{\textregistered}}$ 7, x86-64 (serial and parallel)
  \item Microsoft Windows 10${}^{\mbox{\textregistered}}$, x86-64 (serial)
  \item Apple${}^{\mbox{\textregistered}}$ macOS, x86-64 (serial and parallel)
\end{XyceItemize}


\subsection*{Build Support}
Though not certified platforms, \Xyce{} has been known to run on the following
systems.
\begin{XyceItemize}
  \item FreeBSD 12.X on Intel x86-64 and AMD64 architectures (serial
    and parallel)
  \item Distributions of Linux other than Red Hat Enterprise Linux 6
  \item Microsoft Windows under Cygwin and MinGW
\end{XyceItemize}


\section{\Xyce{} Release \XyceVersionVar{} Documentation}
The following \Xyce{} documentation is available on the \Xyce{} website in pdf
form.
\begin{XyceItemize}
  \item \Xyce{} Version \XyceVersionVar{} Release Notes (this document)
  \item \Xyce{} Users' Guide, Version \XyceVersionVar{}
  \item \Xyce{} Reference Guide, Version \XyceVersionVar{}
  \item \Xyce{} Mathematical Formulation
  \item Power Grid Modeling with \Xyce{}
  \item Application Note: Coupled Simulation with the \Xyce{} General
    External Interface
  \item Application Note: Mixed Signal Simulation with \Xyce{} 7.2
\end{XyceItemize}
Also included at the \Xyce{} website as web pages are the following.
\begin{XyceItemize}
  \item Frequently Asked Questions
  \item Building Guide (instructions for building \Xyce{} from the source code)
  \item Running the \Xyce{} Regression Test Suite
  \item \Xyce{}/ADMS Users' Guide
  \item Tutorial:  Adding a new compact model to \Xyce{}
\end{XyceItemize}


\section{External User Resources}
\begin{itemize}
  \item Website: {\color{XyceDeepRed}\url{http://xyce.sandia.gov}}
  \item Google Groups discussion forum:
    {\color{XyceDeepRed}\url{https://groups.google.com/forum/#!forum/xyce-users}}
  \item Email support:
    {\color{XyceDeepRed}\href{mailto:xyce@sandia.gov}{xyce@sandia.gov}}
  \item Address:
    \begin{quote}
            Electrical Models and Simulation Dept.\\
            Sandia National Laboratories\\
            P.O. Box 5800, M.S. 1177\\
            Albuquerque, NM 87185-1177 \\
    \end{quote}
\end{itemize}

\vspace*{\fill}
\noindent
Sandia National Laboratories is a multimission laboratory managed and
operated by National Technology and Engineering Solutions of Sandia,
LLC, a wholly owned subsidiary of Honeywell International, Inc., for
the U.S. Department of Energy's National Nuclear Security
Administration under contract DE-NA0003525.

\vspace{0.2in}
\noindent
\parbox{\textwidth}
{
    \includegraphics[height=0.6in]{DOEbwlogo}
    \hfill
    \includegraphics[height=0.5in]{NNSA_logo}
    \hfill
    \includegraphics[height=0.5in]{snllineblk}
}
\vspace{0.1in}

\parbox{\textwidth}{\hfill SAND2023-07328 O}
\end{document}

